\chapter{Análisis y discusión de resultados}
\label{capitulo:resultados}
En este capítulo se presentan y analizan los resultados de las ejecuciones de entrenamiento y evaluación de los algoritmos estudiados.
Con el objetivo de distinguir con claridad las etapas metodológicas del entrenamiento de los agentes, la sección se organiza en tres partes.

En primer lugar, se reportan los resultados del proceso de ajuste de hiperparámetros. En segundo lugar, se evalúan los modelos finales seleccionados en los escenarios planteados, desde redes sintéticas hasta el escenario ubicado en Ciudad de Mendoza. Finalmente, se desarrolla una discusión integradora de los hallazgos, sus implicancias y las principales limitaciones del estudio.


\section{Resultados del ajuste de hiperparámetros} 

En esta sección se resume el proceso de ajuste de hiperparámetros y se justifica la configuración utilizada en la evaluación posterior de los modelos.

\subsection{Diseño del proceso de ajuste}

\subsection{Resultados del ajuste y análisis de sensibilidad} 

\subsection{Configuración seleccionada para la evaluación final}  

\section{Resultados de los modelos finales} 

Esta sección reúne los resultados obtenidos con los modelos finales una vez fijada la configuración experimental.
La exposición se organiza según el tipo de escenario analizado. 

\subsection{Criterios de evaluación y métricas de comparación} 

\subsection{Desempeño en escenarios sintéticos: \texorpdfstring{$3\times1$, $3\times3$, $4\times4$ heterogéneo}{3x1, 3x3, 4x4 heterogéneo}}

\subsection{Desempeño en la red de carreteras de la Ciudad de Mendoza} 

Antes de evaluar los controladores, se verificó la operación del escenario canónico v14 de la Ciudad de Mendoza en los dos horizontes definidos en la sección \ref{subsubsec:mendoza_validacion_operacional}. La tabla \ref{tab:metricas_validacion_mendoza} presenta los conteos y las tasas obtenidas con el controlador de tiempo fijo; estos valores caracterizan la condición inicial del escenario y no el desempeño de las políticas aprendidas.

\begin{table}[H]
\centering
\caption{Métricas operacionales del escenario canónico v14 de la Ciudad de Mendoza bajo control de tiempo fijo.}
\label{tab:metricas_validacion_mendoza}
\setlength{\tabcolsep}{4pt}
\small
\begin{adjustbox}{max width=\textwidth}
\begin{tabular}{|l|c|c|c|c|c|c|c|}
\hline
        \textbf{Ventana temporal} & \textbf{Cargados} & \textbf{Insertados} & \textbf{Completados} & \textbf{\acrshort{acr-cr}} & \textbf{Inserción} & \textbf{Completitud} & \textit{teleport} \\ \hline
2 h ($06:30\text{--}08:30$) & 30.629 & 23.735 & 22.702 & 0,7412 & 0,7749 & 0,9565 & 0 \\ \hline
12 h ($06:00\text{--}18:00$) & 180.163 & 143.819 & 143.022 & 0,7938 & 0,7983 & 0,9945 & 0 \\ \hline
\end{tabular}
\end{adjustbox}
\end{table}

\section{Discusión e interpretación de resultados}

En esta sección se integran los resultados presentados y se discuten sus principales implicancias en relación con los objetivos del trabajo.
